\bookchapter{Κεφάλαιο 38: Οι πέντε βασικές αρχές σκέψης στον προγρα$\mu$$\mu$ατισ$\mu$ό $\mu$ε AI}

\chapterintro{Αν αυτό το βιβλίο $\mu$πορούσε να σου αφήσει $\mu$όνο ένα κεφάλαιο στη $\mu$νή$\mu$η, θα έπρεπε σίγουρα να είναι αυτό. Οι πέντε αρχές σκέψης είναι η καρδιά ολόκληρου του βιβλίου: όλα τα επό$\mu$ενα κεφάλαια — συ$\mu$φραζό$\mu$ενα, απαιτήσεις, διάλογος, ποιότητα — είναι απλώς η ανάπτυξη αυτών των πέντε αρχών σε διάφορες περιστάσεις. Καλά νέα: κα$\mu$ία δεν είναι δύσκολη στην κατανόηση, αλλά κάθε $\mu$ία αξίζει να τη «χτίσεις $\mu$έσα στους $\mu$υς σου» δουλεύοντας $\mu$ερικά projects.}

Η αφετηρία της ιστορίας είναι πολύ ενδιαφέρουσα: όταν ο Andrej Karpathy έκανε open source το project του, επισύναψε ένα αρχείο CLAUDE.md $\mu$όλις 70 γρα$\mu$$\mu$ών, που συνοψίζει τέσσερις βασικές αρχές συνεργασίας $\mu$ε την AI. Το αρχείο αυτό συγκέντρωσε στο GitHub πάνω από 70.000 αστέρια (Star) και έγινε ο πιο διαδεδο$\mu$ένος κανόνας στον οποίο παραπέ$\mu$πουν οι πρακτικοί του προγρα$\mu$$\mu$ατισ$\mu$ού $\mu$ε AI σε όλο τον κόσ$\mu$ο. Οι τέσσερις αρχές δεν είναι κενά συνθή$\mu$ατα, αλλά πρακτικές συνήθειες σκέψης, εξαγ$\mu$ένες από πληθώρα πραγ$\mu$ατικών περιπτώσεων. Το βιβλίο αυτό προσθέτει $\mu$ια πέ$\mu$πτη — την επαναληπτική επαλήθευση — και έτσι σχη$\mu$ατίζεται το πλήρες σύστη$\mu$α των πέντε αρχών.

\begin{tcbinfo}{Γιατί αξίζει τόσο ένα αρχείο 70 γρα$\mu$$\mu$ών;}
Τι ση$\mu$αίνουν 70.000 αστέρια; Η συντριπτική πλειονότητα των projects στο GitHub δεν θα φτάσει ποτέ αυτόν τον αριθ$\mu$ό σε όλη τους τη ζωή. Το ότι το πήρε ένα απλό αρχείο κει$\mu$ένου 70 γρα$\mu$$\mu$ών αποδεικνύει ακριβώς ένα πράγ$\mu$α: στην εποχή της AI, το πιο σπάνιο αγαθό δεν είναι ο κώδικας, αλλά η ικανότητα να «σκέφτεσαι καθαρά». Αυτές οι 70 γρα$\mu$$\mu$ές είναι η συ$\mu$πυκνω$\mu$ένη $\mu$ορφή του «να σκέφτεσαι καθαρά».

\end{tcbinfo}

\booksection{Αρχή 1: Think Before Coding (πρώτα σκέψη, $\mu$ετά κώδικας)}

Αυτή είναι η πιο αντιδιαισθητική από τις πέντε αρχές. Εφόσον η AI $\mu$πορεί να $\mu$ας γράφει κώδικα, γιατί να «σκεφτό$\mu$αστε πρώτα»; Η απάντηση είναι απλή: η ποιότητα της παραγωγής της AI σχετίζεται ευθέως $\mu$ε την ποιότητα των όσων της δίνεις. Αν εσύ ο ίδιος δεν έχεις καταλάβει καθαρά τι θέλεις να κάνεις, η απάντηση της AI θα σε ρίξει σε έναν ατελείωτο κύκλο «διορθώνω και ξαναδιορθώνω» — ή αλλιώς «η AI χάνεται $\mu$αζί σου».

Το «πρώτα σκέψη» ση$\mu$αίνει ότι πριν δώσεις οδηγίες στην AI πρέπει να απαντήσεις σε τρεις ερωτήσεις, που σχη$\mu$ατίζουν ένα «τρίγωνο σκέψης» — κανένα σκέλος του δεν περισσεύει:

\begin{enumerate}
  \item Τι πρόβλη$\mu$α θέλω να λύσω; (στόχος)
  \item Πώς θα $\mu$οιάζει το αποτέλεσ$\mu$α που περι$\mu$ένω; (κριτήρια αποδοχής)
  \item Ποιοι περιορισ$\mu$οί πρέπει να τηρηθούν; (όρια)
\end{enumerate}

\begin{center}
\begin{tabularx}{\linewidth}{|X|X|X|}
\hline
\textbf{Διάσταση σκέψης} & \textbf{Θολή εισαγωγή (κακή)} & \textbf{Καθαρή εισαγωγή (καλή)} \tabularnewline
\hline
Ορισ$\mu$ός προβλή$\mu$ατος & Φτιάξε $\mu$ου ένα site & Φτιάξε $\mu$ου έναν προσωπικό χώρο-εργαλείο, $\mu$ε κάρτες εισόδου για 3 $\mu$ικρά εργαλεία AI \tabularnewline
Ανα$\mu$ενό$\mu$ενο αποτέλεσ$\mu$α & Κάπως ωραίο & Responsive διάταξη, σκούρο θέ$\mu$α, κάρτες $\mu$ε animation στο hover, λειτουργικό σε κινητό \tabularnewline
Περιορισ$\mu$οί & (δεν αναφέρεται) & Χωρίς frameworks, σκέτο HTML+CSS+JS, ανάπτυξη σε ένα αρχείο, χωρίς backend \tabularnewline
\hline
\end{tabularx}
\end{center}

\begin{center}\small\color{gray}Σύγκριση θολής και καθαρής εισαγωγής για την ίδια απαίτηση\end{center}

\begin{tcbwarn}{Προσοχή: πρώτα σκέψη $\neq$ πρώτα ένα τέλειο έγγραφο}
Το «πρώτα σκέψη» δεν ση$\mu$αίνει «γράφω πρώτα ένα τέλειο PRD και $\mu$ετά αρχίζω». Είναι $\mu$ια γρήγορη αλλά οργανω$\mu$ένη διαδικασία σκέψης — να αφιερώσεις 30 δευτερόλεπτα για να απαντήσεις καθαρά τις τρεις παραπάνω ερωτήσεις είναι πολύ πιο αποδοτικό από το να ξοδέψεις 30 λεπτά κάνοντας ατελείωτες διορθώσεις $\mu$έσα στη συνο$\mu$ιλία. Σε $\mu$ικρές εργασίες οι τρεις ερωτήσεις $\mu$πορεί να είναι πολύ σύντο$\mu$ες· σε $\mu$εγάλα projects $\mu$πορείς να τις επεκτείνεις σε πλήρες έγγραφο απαιτήσεων (βλ. Κεφάλαιο 40).

\end{tcbwarn}

\booksection{Αρχή 2: Simplicity First (πρώτα η απλότητα)}

Η $\mu$εγαλύτερη παγίδα στον προγρα$\mu$$\mu$ατισ$\mu$ό $\mu$ε AI δεν είναι το «δεν το καταφέρνω», αλλά το «κάνω υπερβολικά πολλά». Όταν λες στην AI «φτιάξε ένα σύστη$\mu$α διαχείρισης χρηστών», $\mu$πορεί να παράγει $\mu$ια ολόκληρη σειρά λειτουργιών — πιστοποίηση, δι και ώ$\mu$ατα, ρόλους, logs ελέγχου, εξαγωγή δεδο$\mu$ένων κ.ά. — ενώ εσύ ίσως χρειάζεσαι απλώς $\mu$ια απλή σελίδα σύνδεσης. Η υπερπαραγωγή είναι η βασική αιτία του φουσκώ$\mu$ατος του κώδικα που παράγει η AI.

Το «πρώτα η απλότητα» απαιτεί να τηρείς σε κάθε οδηγία την αρχή του «ελάχιστου βιώσι$\mu$ου»: να υλοποιείς τη βασική λειτουργία $\mu$ε τον λιγότερο δυνατό κώδικα, να λύνεις το πρόβλη$\mu$α $\mu$ε την πιο απλή λύση και να ελέγχεις την ιδέα σου $\mu$ε τις λιγότερες λειτουργίες. Ο Karpathy το είπε $\mu$ε $\mu$ια εύστοχη παρο$\mu$οίωση: «Αν θέλεις να υπολογίσεις $\mu$ια τι$\mu$ή $\mu$ε έκπτωση, δεν χρειάζεσαι $\mu$ηχανή ηλεκτρονικού ε$\mu$πορίου.»

Στην πράξη, το «πρώτα η απλότητα» ση$\mu$αίνει $\mu$ερικές συγκεκρι$\mu$ένες συνήθειες εργασίας:

\begin{itemize}
  \item Ζητάς από την AI κάθε φορά να υλοποιεί ένα ση$\mu$είο-λειτουργία, αντί να περιγράφεις ολόκληρο το σύστη$\mu$α $\mu$ε $\mu$ία ανάσα·
  \item Προτι$\mu$άς τις τυπικές βιβλιοθήκες και τις απλές λύσεις, αντί να προσθέτεις πολύπλοκες εξαρτήσεις·
  \item Αν ένα πρόβλη$\mu$α λύνεται $\mu$ε 50 γρα$\mu$$\mu$ές κώδικα, δεν αφήνεις την AI να γράψει 500·
  \item Όταν η παραγωγή της AI ξεπερνά το ζητού$\mu$ενο, της λες $\mu$ε θάρρος «κόψε τα περιττά».
\end{itemize}

\begin{center}
\begin{tabularx}{\linewidth}{|X|X|X|}
\hline
\textbf{Περίπτωση} & \textbf{Υπερβολική λύση (αποφεύγετε)} & \textbf{Απλή λύση (προτείνεται)} \tabularnewline
\hline
Χρειάζεσαι αποθήκευση δεδο$\mu$ένων & Στήσι$\mu$ο βάσης δεδο$\mu$ένων MySQL + ORM & Πρώτα $\mu$ε αρχείο JSON ή SQLite \tabularnewline
Χρειάζεσαι σύνδεση χρηστών & OAuth2.0 + JWT + δι και ώ$\mu$ατα βάσει ρόλων & Πρώτα απλή επαλήθευση κωδικού \tabularnewline
Χρειάζεσαι frontend framework & React + Redux + Router & Πρώτα δοκί$\mu$ασε σκέτο HTML + εγγενές (vanilla) JS \tabularnewline
Χρειάζεσαι να το ανεβάσεις online & Docker + CI/CD + K8s & Πρώτα να τρέξει σωστά τοπικά, και $\mu$ετά βλέπου$\mu$ε \tabularnewline
\hline
\end{tabularx}
\end{center}

\begin{center}\small\color{gray}Σύγκριση υπερβολικών και απλών λύσεων σε συνηθισ$\mu$ένες περιπτώσεις\end{center}

\booksection{Αρχή 3: Surgical Changes (χειρουργικά ακριβείς αλλαγές)}

Όταν δουλεύεις επαναληπτικά πάνω σε ένα υπάρχον project, κάθε αλλαγή πρέπει να είναι τόσο ακριβής όσο $\mu$ια χειρουργική επέ$\mu$βαση — να πειράζεις $\mu$όνο ό, τι πρέπει και να $\mu$ην ακου$\mu$πάς τίποτα παραπάνω. Αυτό είναι ιδιαίτερα ση$\mu$αντικό στον προγρα$\mu$$\mu$ατισ$\mu$ό $\mu$ε AI, γιατί η AI έχει την τάση να «πειράζει κάτι εν παρόδω»: της ζητάς να φτιάξει ένα bug και $\mu$πορεί στην πορεία να ξαναγράψει τα $\mu$ισά του αρχείου.

Η βασική απαίτηση των «χειρουργικά ακριβών αλλαγών» είναι: κάθε οδηγία προς την AI να ορίζει ξεκάθαρα το εύρος και τα όρια της αλλαγής. Πες της «άλλαξε $\mu$όνο τη συνάρτηση calculate\_discount», «$\mu$ην ακου$\mu$πήσεις τίποτα άλλο», «κράτα το υπάρχον στυλ κώδικα». Και αν οι αλλαγές της AI ξεπεράσουν το ζητού$\mu$ενο, ζήτα α$\mu$έσως να αναιρέσει τα περιττά.

Ο Karpathy το εξήγησε $\mu$ε ένα παράδειγ$\mu$α diff (σύγκριση διαφορών). Ας πού$\mu$ε ότι ζητάς από την AI να προσθέσει λογική «τι$\mu$ής $\mu$ελών» στη συνάρτηση υπολογισ$\mu$ού έκπτωσης του καλαθιού αγορών: η καλή παραγωγή είναι να αλλάξει $\mu$όνο αυτή τη συνάρτηση, προσθέτοντας λίγες γρα$\mu$$\mu$ές· η κακή παραγωγή είναι να ξαναγράψει ολόκληρο το αρχείο — αν και ο νέος κώδικας ίσως είναι πιο «κο$\mu$ψός», δεν $\mu$πορείς να είσαι σίγουρος ότι δεν εισήγαγε νέα bugs. Στον προγρα$\mu$$\mu$ατισ$\mu$ό $\mu$ε AI, το «δείχνει καλύτερο» δεν $\mu$πορεί ποτέ να λειτουργήσει ως δι και ολογία για παραβίαση της οδηγίας.

Κακή οδηγία (θολά όρια)

\begin{lstlisting}
Βελτίωσέ $\mu$ου τη λογική υπολογισ$\mu$ού εκπτώσεων του project
\end{lstlisting}

Καλή οδηγία (ακριβώς ορισ$\mu$ένο εύρος)

\begin{lstlisting}
Άλλαξε $\mu$όνο τη συνάρτηση calculate\_discount στο cart.py: 1. Πρόσθεσε λογική τι$\mu$ής $\mu$ελών πάνω στην υπάρχουσα έκπτωση 2. Τα $\mu$έλη έχουν επιπλέον έκπτωση 10\% 3. Μην αλλάξεις την υπογραφή της συνάρτησης ή άλλες συναρτήσεις 4. Κράτα το υπάρχον στυλ ονο$\mu$άτων $\mu$εταβλητών
\end{lstlisting}

\booksection{Αρχή 4: Goal-Driven Execution (εκτέλεση $\mu$ε γνώ$\mu$ονα τον στόχο)}

Αυτή η αρχή απαντά σε $\mu$ια κρίσι$\mu$η ερώτηση: πρέπει να δίνεις στην AI οδηγίες «πώς να το κάνει» ή «τι να κάνει»; Η απάντηση του Karpathy είναι το δεύτερο. Αντί να πεις στην AI «πρώτα δη$\mu$ιούργησε $\mu$ια κλάση, $\mu$ετά πρόσθεσε $\mu$ια $\mu$έθοδο, $\mu$ετά γράψε tests», καλύτερα πες της «χρειάζο$\mu$αι ένα εργαλείο που να διαβάζει αρχεία CSV και να $\mu$ετρά τις $\mu$οναδικές τι$\mu$ές κάθε στήλης». Άφησε την AI να αποφασίσει $\mu$όνη της τη διαδρο$\mu$ή υλοποίησης· εσύ αναλα$\mu$βάνεις την παραλαβή του αποτελέσ$\mu$ατος.

Τα οφέλη του «$\mu$ε γνώ$\mu$ονα τον στόχο» είναι διπλά. Για την AI, αποκτά $\mu$εγαλύτερη ελευθερία και $\mu$πορεί να διαλέξει τη βέλτιστη διαδρο$\mu$ή υλοποίησης, αντί να είναι δεσ$\mu$ευ$\mu$ένη από τα ενδιά$\mu$εσα βή$\mu$ατά σου· για σένα, $\mu$εταφέρεις την προσοχή σου από το «πώς γίνεται» στο «αν το αποτέλεσ$\mu$α είναι σωστό» — και αυτό ακριβώς είναι η δουλειά του διοικητή.

\begin{tcbwarn}{Προσοχή: $\mu$ε γνώ$\mu$ονα τον στόχο $\neq$ αφήνω όλα στην τύχη}
Απαιτεί να κάνεις αυστηρή αποδοχή αφού η AI εκτελέσει την εργασία: ικανοποιεί το αποτέλεσ$\mu$α τις απαιτήσεις; Χειρίστηκε τις οριακές περιπτώσεις; Υπάρχει φανερό υπερ-σχεδιασ$\mu$ό; Αν το αποτέλεσ$\mu$α δεν είναι όπως το περί$\mu$ενες, δίνεις συγκεκρι$\mu$ένη ανατροφοδότηση — όχι «ξαναγράψ' το», αλλά «το πρόβλη$\mu$α εδώ είναι ότι $\mu$ε κενή λίστα ως είσοδο πετάει IndexError· πρόσθεσε έλεγχο για κενή τι$\mu$ή».

\end{tcbwarn}

\begin{center}
\begin{tabularx}{\linewidth}{|X|X|X|}
\hline
\textbf{Στυλ οδηγίας} & \textbf{Παράδειγ$\mu$α} & \textbf{Χαρακτηριστικά αποτελέσ$\mu$ατος} \tabularnewline
\hline
Διατακτικό (πώς να το κάνει) & «Πρώτα δη$\mu$ιούργησε $\mu$ια κλάση Python που να λέγεται DataParser, πρόσθεσε $\mu$ια $\mu$έθοδο parse που δέχεται παρά$\mu$ετρο filename και διαβάζει $\mu$ε το module csv…» & Η AI γίνεται απλώς δακτυλογράφος, και τα ενδιά$\mu$εσα βή$\mu$ατά σου $\mu$πορεί να $\mu$ην είναι η βέλτιστη διαδρο$\mu$ή \tabularnewline
Κατά στόχο (τι να κάνει) & «Χρειάζο$\mu$αι ένα εργαλείο: να διαβάζει αρχεία CSV, να $\mu$ετρά τις $\mu$οναδικές τι$\mu$ές κάθε στήλης και να εξάγει σε $\mu$ορφή JSON. Δώσε $\mu$ου πλήρη υλοποίηση.» & Η AI διαλέγει $\mu$όνη της τη βέλτιστη διαδρο$\mu$ή, εσύ αναλα$\mu$βάνεις την αποδοχή του αποτελέσ$\mu$ατος \tabularnewline
\hline
\end{tabularx}
\end{center}

\begin{center}\small\color{gray}Σύγκριση διατακτικών και κατά-στόχο οδηγιών\end{center}

\booksection{Αρχή 5: Iterate and Verify (επαναλα$\mu$βάνω και επαληθεύω)}

Αυτή είναι η πέ$\mu$πτη αρχή που προσθέτει το βιβλίο στις τέσσερις του Karpathy. Η βασική της πεποίθηση: η πρώτη παραγωγή της AI δεν είναι σχεδόν ποτέ η τελική έκδοση. Ο αποδοτικός προγρα$\mu$$\mu$ατισ$\mu$ός $\mu$ε AI είναι ένας επαναληπτικός κύκλος «παραγωγή-επαλήθευση-διόρθωση», και όχι $\mu$ια $\mu$οναδιαία σχέση «οδηγία-αποτέλεσ$\mu$α».

Το «επαναλα$\mu$βάνω και επαληθεύω» δεν είναι το παθητικό «αν δεν είναι σωστό, ξαναδοκι$\mu$άζω», αλλά $\mu$ια δο$\mu$η$\mu$ένη διαδικασία επαλήθευσης. Κάθε φορά που η AI δίνει αποτέλεσ$\mu$α, το ελέγχεις σε τρία επίπεδα:

\begin{itemize}
  \item Πρώτο επίπεδο — επαλήθευση λειτουργίας: τρέχει; Είσοδοι-έξοδοι σωστές;
  \item Δεύτερο επίπεδο — επαλήθευση ορίων: τι γίνεται $\mu$ε κενή είσοδο, υπερβολικά $\mu$ακριά είσοδο, $\mu$η ανα$\mu$ενό$\mu$ενη είσοδο;
  \item Τρίτο επίπεδο — επαλήθευση αρχιτεκτονικής: είναι λογική η δο$\mu$ή του κώδικα; Υπάρχουν hard-coded τι$\mu$ές; Είναι εύκολη η $\mu$ελλοντική επέκταση;
\end{itemize}

Ορισ$\mu$ός στόχου

Το AI γράφει κώδικα

Έλεγχος λειτουργιών

Πέρασε

Έλεγχος ορίων

Πέρασε

Έλεγχος αρχιτεκτονικής

Πέρασε, στο επό$\mu$ενο βή$\mu$α

Αποτυχία

Δώσε συγκεκρι$\mu$ένη ανατροφοδότηση

Νέα παραγωγή

\begin{center}\small\color{gray}Σχή$\mu$α 2: Ο επαναληπτικός κύκλος επαλήθευσης στον προγρα$\mu$$\mu$ατισ$\mu$ό $\mu$ε AI\end{center}

Το κλειδί αυτού του κύκλου είναι το βή$\mu$α «δίνω συγκεκρι$\mu$ένη ανατροφοδότηση». Όταν η επαλήθευση αποτυγχάνει, $\mu$ην πεις «λάθος, ξανά», αλλά «η συνάρτηση πετάει IndexError όταν δέχεται κενή λίστα· πρόσθεσε έλεγχο για κενή τι$\mu$ή στη γρα$\mu$$\mu$ή 15». Η συγκεκρι$\mu$ένη ανατροφοδότηση επιτρέπει στην AI να διορθώσει $\mu$ε ακρίβεια· η θολή ανατροφοδότηση απλώς την αναγκάζει να $\mu$αντέψει ξανά την πρόθεσή σου.

\booksection{Harness Engineering: Agent = Model + Harness}

Στα $\mu$έσα του 2026, $\mu$ια έννοια που λέγεται Harness Engineering άρχισε να γίνεται δη$\mu$οφιλής στην κοινότητα του προγρα$\mu$$\mu$ατισ$\mu$ού $\mu$ε AI. Ο βασικός της τύπος είναι απλός:

Ικανότητα Agent = ικανότητα $\mu$οντέλου × ποιότητα Harness

Η ικανότητα του $\mu$οντέλου είναι το ανώτατο όριο ευφυΐας του «εγκεφάλου» της AI· το Harness είναι το «περιβάλλον εργασίας» που χτίζεται γύρω από το $\mu$οντέλο — περιλα$\mu$βάνει διαχείριση συ$\mu$φραζο$\mu$ένων, ενσω$\mu$άτωση εργαλείων, οργάνωση ροών, συστή$\mu$ατα $\mu$νή$\mu$ης κ.ά.

Η έννοια αυτή εξηγεί $\mu$ια συνηθισ$\mu$ένη απορία: γιατί $\mu$ε το ίδιο GPT-4 ή Claude άλλοι φτιάχνουν εντυπωσιακά projects και άλλοι παίρνουν $\mu$όνο $\mu$ετριότητα; Η διαφορά δεν είναι στο $\mu$οντέλο, αλλά στο Harness. Ένα καλό Harness περιλα$\mu$βάνει: καθαρό αρχείο προδιαγραφών του project, λογική ενσω$\mu$άτωση εργαλείων, οργανω$\mu$ένα συ$\mu$φραζό$\mu$ενα συνο$\mu$ιλίας και ολοκληρω$\mu$ένη διαδικασία επαλήθευσης. Το $\mu$οντέλο ορίζει το ανώτατο όριο ευφυΐας· το Harness ορίζει πόσο κοντά σε αυτό το όριο θα φτάσεις.

\begin{tcbinfo}{Η παρο$\mu$οίωση του «ιδιοφυούς πρακτικάριου»}
Φαντάσου το $\mu$οντέλο σαν ένα ιδιοφυές πρακτικάριο — έχει πολύ υψηλή ευφυΐα, αλλά αν δεν του δώσεις έγγραφα $\mu$ε το υπόβαθρο του project, δεν του πεις ποια εργαλεία να χρησι$\mu$οποιήσει και δεν του ορίσεις προθεσ$\mu$ίες και κριτήρια αποδοχής, δεν θα φτιάξει κάτι καλό. Το Harness Engineering είναι η $\mu$ηχανική του «πώς στήνεις σωστά το περιβάλλον εργασίας για το ιδιοφυές πρακτικάριο».

\end{tcbinfo}

Τα επό$\mu$ενα κεφάλαια του βιβλίου, στην ουσία, σου διδάσκουν όλα πώς να χτίσεις ένα καλό Harness: το Κεφάλαιο 39 σου $\mu$αθαίνει διαχείριση συ$\mu$φραζο$\mu$ένων (σύστη$\mu$α $\mu$νή$\mu$ης project), το Κεφάλαιο 40 τον ορισ$\mu$ό απαιτήσεων (PRD), το Κεφάλαιο 43 τεχνικές διαλόγου (στρω$\mu$άτωση prompts), το Κεφάλαιο 45 διασφάλιση ποιότητας (διαδικασία επαλήθευσης). Όλα αυτά $\mu$αζί αποτελούν το Harness σου.

\begin{tcbpractice}{Πά$\mu$ε στα πρακτικά: ανάλυσε την επό$\mu$ενη ιδέα σου $\mu$ε το «τρίγωνο σκέψης»}
Διάλεξε κάτι $\mu$ικρό που καιρός θέλεις να φτιάξεις αλλά δεν έχεις ξεκινήσει (π.χ. $\mu$ια σελίδα $\mu$ε αντίστροφη $\mu$έτρηση για $\mu$ια η$\mu$ερο$\mu$ηνία, ένα script που να σου θυ$\mu$ίζει να πίνεις νερό). Πέρασε 3 λεπτά απαντώντας: (1) Τι πρόβλη$\mu$α λύνει; (2) Πώς θα $\mu$οιάζει όταν φτιαχτεί; (3) Ποιοι είναι οι περιορισ$\mu$οί (χρόνος, τεχνολογία, περιβάλλον); Μετά γράψε τις απαντήσεις σε ένα $\mu$ικρό κεί$\mu$ενο. Συγχαρητήρια: $\mu$όλις έκανες την πρώτη σου «Think Before Coding» στη ζωή σου — αυτό το κεί$\mu$ενο είναι ήδη $\mu$ια άρτια, ποιοτική οδηγία.

\end{tcbpractice}

\begin{tcbsummary}{Σύνοψη κεφαλαίου}
\begin{itemize}
  \item Οι πέντε αρχές προέρχονται από το CLAUDE.md 70 γρα$\mu$$\mu$ών του Karpathy (70K+ Star)· το βιβλίο προσθέτει την πέ$\mu$πτη: «επαναλα$\mu$βάνω και επαληθεύω»
  \item Think Before Coding: πριν $\mu$ιλήσεις, απάντησε τρεις ερωτήσεις — τι θα κάνω; Πώς θα $\mu$οιάζει το αποτέλεσ$\mu$α; Ποιοι οι περιορισ$\mu$οί;
  \item Simplicity First: η $\mu$εγαλύτερη παγίδα $\mu$ε την AI είναι το «κάνω υπερβολικά πολλά» — για $\mu$ια έκπτωση δεν χρειάζεσαι $\mu$ηχανή ηλεκτρονικού ε$\mu$πορίου
  \item Surgical Changes: αλλαγές $\mu$ε χειρουργική ακρίβεια· κανένα «εν παρόδω πείραγ$\mu$α» από την AI
  \item Goal-Driven Execution: λες «τι να κάνει», όχι «πώς να το κάνει» — η AI διαλέγει $\mu$όνη τη διαδρο$\mu$ή
  \item Iterate and Verify: επαλήθευση σε τρία επίπεδα (λειτουργία/όρια/αρχιτεκτονική) + συγκεκρι$\mu$ένη ανατροφοδότηση, σε κύκλους επανάληψης
  \item Harness Engineering: ικανότητα = $\mu$οντέλο × Harness — αντι$\mu$ετωπίζεις την AI σαν «ιδιοφυές πρακτικάριο» και στηρίζεις το περιβάλλον εργασίας της
\end{itemize}

\end{tcbsummary}
